\newpage
\section{Construcciones de Fibrados Vectoriales a Partir de Otros}

\cite[\S 3]{ClasesCaracteristicas}.

\subsection{Restricción de fibrados}

\subsection{Fibrados Inducidos}

Sea $\xi$ un fibrado sobre $B$, y sea $B_{1}$ un espacio topológico arbitrario. Dado un mapa $f \colon B_{1} \to B$ cualquiera, podemos construir el fibrado inducido $f^{*} \xi$ sobre $B_{1}$. El espacio total $E_{1}$ de $f^{*} \xi$ es el subconjunto $E_{1} \subset B_{1} \times E$, que consiste de todos los pares $(b,e)$, con
\[ f(b) = \pi(e). \]
El mapa de proyección $\pi_{1} \colon E_{1} \to B_{1}$ se define por $\pi_{1}(b,e) = b$. Luego, tenemos un diagrama conmutativo
\[ \begin{tikzcd}
E_1 \arrow[r, "\hat{f}", xshift=-2pt] \arrow[d, "\pi_1"', xshift=-4pt]      & E \arrow[d, "\pi", xshift=-4pt] \\
B_1 \arrow[r, "f"{xshift=-4pt}, shorten <=-4pt, shorten >=4pt, xshift=-2pt] & B
\end{tikzcd} \]
donde $\hat{f}(b,e) = e$. La estructura de espacio vectorial en $\pi_{1}^{-1}(b)$ está definida por
\[ t_{1}(b,e_{1}) + t_{2}(b,e_{2}) = (b, t_{1}e_{1} + t_{2}e_{2}). \]
Notemos que, con esta estructura, se da el caso que $e_{1},e_{2} \in E$ pertenecen a la misma fibra en $\xi$, dado que el diagrama anterior conmuta. Luego, no es dificil ver que $\hat{f}$ envía el espacio vectorial $F_{b}(f^{*}\xi)$ al espacio vectorial $F_{f(b)}(\xi)$ de manera isomorfa. \par
Para asegurarnos de que esta construcción nos está dando, en efecto, un nuevo fibrado vectorial, debemos verificar que se cumple la condición de trivialización local. Es decir, si $(U,h)$ es un sistema de coordenadas locales para $\xi$, definimos $U_{1} = f^{-1}(U)$ y
\[ h_{1} \colon U_{1} \times \R^{n} \to \pi_{1}^{-1}(U_{1}), \]
de tal manera que $h_{1}(b,x) = (b, h(f(b), x))$.
%Aún no estoy seguro de si quiero dejar esto en un lemma o continuar explicando de manera natural. Supongo que esto va a depender de si mi demostración se extiende mucho o no.
\begin{lemma}
	$(U_{1}, h_{1})$ es un sistema de coordenadas locales para $f^{*}\xi$.
\end{lemma}
\begin{proof}
	...
\end{proof}

Es importante notar que si $\xi$ es trivial, entonces $f^{*}\xi$ también lo será. \par
El anterior diagrama conmutativo sugiere el siguiente concepto que, a priori, es más general. El lemma que vendrá luego dejará en claro por qué esto es sólo \textit{a priori}.

\begin{definition}
	Sean $\xi$ y $\eta$ fibrados vectoriales. Un \textbf{mapa de fibrados} desde $\xi$ a $\eta$ es una función continua
	\[ g \colon E(\xi) \to E(\eta), \]
	tal que envía cada espacio vectorial $F_{b}(\xi)$ hacia un espacio vectorial $F_{b'}(\eta)$ de manera isomorfa.
\end{definition}

Definiendo $\bar{g}(b) = b'$, es claro que la función resultante
\[ \bar{g} \colon B(\xi) \to B(\eta), \]
es continua.

\begin{lemma}
	Si $g \colon E(\xi) \to E(\eta)$ es un mapa de fibrados, con $\bar{g}$ su mapa correspondiente entre espacios bases, entonces $\eta$ es isomorfo al fibrado inducido $\bar{g}^{*}\xi$.
\end{lemma}
\begin{proof}
	...
\end{proof}

\subsection{Productos Cartesianos}

Sean $\xi_{1},\xi_{2}$ fibrados con mapas de proyección $\pi_{i} \colon E_{i} \to B_{i}$, $i=1,2$, respectivamente. El \textbf{producto cartesiano} $\xi_{1} \times \xi_{2}$ está definido como el fibrado con mapa de proyección
\[ \pi_{1} \times \pi_{2} \colon E_{1} \times E_{2} \to B_{1} \times B_{2}, \]
donde a cada fibra
\[ (\pi_{1} \times \pi_{2})^{-1}(b_{1},b_{2}) = F_{b_{1}}(\xi_{1}) \times F_{b_{2}}(\xi_{2}) \]
le es dada la estructura de espacio vectorial obvia.

% Acá, debería considerar una de dos, o demostrar la afirmación de ejemplo, o simplemente quitarlo.
\begin{example}
	Dadas dos variedades diferenciables $M_{1},M_{2}$, consideremos sus respectivos fibrados tangentes $\tau_{M_{1}}, \tau_{M_{2}}$. Tomando el producto cartesiano, obtenemos el fibrado $\tau_{M_{1}} \times \tau_{M_{2}}$. Luego, si consideramos el fibrado tangente $\tau_{M}$ de la variedad diferenciable $M = M_{1} \times M_{2}$, entonces es posible demostrar que este fibrado es isomorfo a $\tau_{M_{1}} \times \tau_{M_{2}}$.
\end{example}

\subsection{Sumas de Whitney}

Consideremos dos fibrados $\xi_{1}, \xi_{2}$ sobre el mismo espacio base $B$. Sea
\[ \Delta \colon B \to B \times B \]
el encaje sobre la diagonal. El fibrado inducido $\Delta^{*}(\xi_{1} \times \xi_{2})$ sobre $B$ es llamado la \textbf{suma de Whitney}, y será denotado por $\xi_{1} \oplus \xi_{2}$. Es posible notar que cada fibra $F_{b}(\xi_{1} \oplus \xi_{2})$ es canónicamente isomorfa a la suma directa $F_{b}(\xi_{1}) \oplus F_{b}(\xi_{2})$.

\begin{definition}
	Consideremos dos fibrados $\xi$ y $\eta$ sobre el mismo espacio base, con $E(\xi) \subset E(\eta)$. Decimos que $\xi$ es un \textbf{sub-fibrado} de $\eta$ si cada fibra $F_{b}(\xi)$ es un subespacio vectorial del correspondiente fibrado $F_{b}(\eta)$. En tal caso escribimos $\xi \subset \eta$.
\end{definition}

\begin{lemma}
	Sean $\xi_{1}$ y $\xi_{2}$ sub-fibrados de $\eta$, tales que cada espacio vectorial $F_{b}(\eta)$ es igual a la suma directa de espacios vectoriales $F_{b}(\xi_{1}) \oplus F_{b}(\xi_{2})$. Entones, $\eta$ es isomorfo a la suma de Whitney $\xi_{1} \oplus \xi_{2}$.
\end{lemma}
\begin{proof}
	...
\end{proof}

\subsection{Complementos Ortogonales}

A partir de este lemma, es natural preguntarse si, dado un sub-fibrado $\xi \subset \eta$, ¿existe otro sub-fibrado tal que $\eta$ se separa en una suma de Whitney de estos sub-fibrados? Para responder a esto, y tal como se sugirió al final de la sección anterior, nos reduciremos a los fibrados que poseen una métrica Euclidiana. \par

Sea $\eta$ un fibrado vectorial Euclidiano. Definimos $F_{b}(\xi^{\perp})$ como el subespacio de $F_{b}(\eta)$ que consiste de los vectores $v$ que cumplen $v \cdot w = 0$ para todo $w \in F_{b}(\xi)$. Sea $E(\xi^{\perp}) \subset E(\eta)$ la unión de los $F_{b}(\xi^{\perp})$.

\begin{definition}
	$\xi^{\perp}$ será llamado el \textbf{complemento ortogonal} de $\xi$ en $\eta$.
\end{definition}

\begin{theorem}
	$E(\xi^{\perp})$ es el espacio total de un sub-fibrado $\xi^{\perp} \subset \eta$. Aún más, $\eta$ es isomorfo a la suma de Whitney $\xi \oplus \xi^{\perp}$.
\end{theorem}
\begin{proof}
	...
\end{proof}

\subsection{El Funtor $\Hom$}

Es posible construir nuevos espacios vectoriales de muchas formas a partir de otros. Si bien nos centraremos principalmente en el funtor $\Hom$, pues nos será de utilidad para ejemplificar el cálculo de calses de Stiefel-Whitney para ciertos fibrados, introduciremos una forma general de construir fibrados a partir de dos fibrados distintos mediante funtores de espacios vectoriales.

\begin{definition}
	Denotemos por $\Vect_{\R}$ a la categoría de espacios vectoriales reales finito dimensionales y de isomorfismos entre ellos. Un \textbf{funtor} (covariante) $T \colon \Vect_{\R} \times \Vect_{\R} \to \Vect_{\R}$ es una operación que asigna
	\begin{itemize}
		\item[1.] a cada par $V,W \in \Vect_{\R}$ de espacios vectoriales, un espacio vectorial $T(V,W) \in \Vect_{\R}$;

		\item[2.] a cada par $f \colon V \to V',\, g \colon W \to W'$ de isomorfismos, un isomorfismo
		\[ T(f,g) \colon T(V,W) \to T(V',W') \]
		tal que

		\item[3.] $T(\id_{V}, \id_{W}) = \id_{T(V,W)}$ y

		\item[4.] $T(f_{1} \circ f_{2}, g_{1} \circ g_{2}) = T(f_{1}, g_{1}) \circ T(f_{2},g_{2})$.
	\end{itemize}
	Tal funtor será llamado \textbf{continuo} si $T(f,g)$ depende continuamente de $f$ y $g$. Esto tiene sentido, pues el conjunto de todos los isomorfismos de un espacio vectorial hacia otro tiene una topología natural.
\end{definition}

Es posible extender esta definición de funtores continuos a más variables $T : \Vect_{\R} \times \cdots \times \Vect_{\R} \to \Vect_{\R}$ de manera similar. Nosotros nos restringimos al caso en que tenemos sólo dos variables. \par
Ahora, sea $T : \Vect_{\R} \times \Vect_{\R} \to \Vect_{\R}$ un funtor continuo, y sean $\xi_{1},\xi_{2}$ fibrados vectoriales sobre un mismo espacio base $B$. El nuevo fibrado vectorial se construye de la siguiente manera. Para cada $b \in B$, sea
\[ F_{b} = T(F_{b}(\xi_{1}), F_{b}(\xi_{2})). \]
Denotamos por $E$ a la unión disjunta de los espacios vectoriales $F_{b}$ y definimos $\pi \colon E \to B$ por $\pi(F_{b}) = b$. \par

Si bien en el contexto de este informe no estamos completamente interesados en la demostración del siguiente teorema, sigue siendo algo importante a mencionar para que la construcción anterior tenga sentido.

\begin{theorem}
	Existe una topología canónica para $E$ de tal manera que $E$ es el espacio total de un fibrado vectorial con proyección $\pi$ y fibras $F_{b}$.
\end{theorem}

\begin{definition}
	Este fibrado será denotado por $T(\xi_{1},\xi_{2})$.
\end{definition}

En particular, recordemos que $\Hom(V,W)$ es el espacio vectorial de todas las transformaciones lineales entre $V$ y $W$. En efecto, $\Hom : \Vect_{\R} \times \Vect_{\R} \to \Vect_{\R}$ es un funtor continuo, por lo que podemos construir nuevos fibrados vectoriales, y lo denotaremos por $\Hom(\xi_{1},\xi_{2})$. Notemos que, en este caso, cada fibra representará el espacio de transformaciones lineales entre las fibras $F_{b}(\xi_{1})$ y $F_{b}(\xi_{2})$. \par

Una herramienta que nos será de utilidad durante en una sección próxima es el siguiente lemma.

\begin{lemma}
	Si un fibrado vectorial $\xi$ posee una métrica Euclidiana, entonces $\xi$ es isomorfo al fibrado $\Hom(\xi, \varepsilon^{1})$.
\end{lemma}
\begin{proof}
	...
\end{proof}
